\documentclass[11pt]{article}

\usepackage[margin=1in]{geometry}
\usepackage{amsmath, amssymb, amsthm}
\usepackage{booktabs}
\usepackage{longtable}
\usepackage{enumitem}
\usepackage{microtype}
\usepackage{parskip}
\usepackage{xcolor}  % needed for blue!50!black color mixing in hyperref setup
\usepackage[round, authoryear]{natbib}  % author-year citations

% XeLaTeX font setup --- uses local Times fonts from the fonts/ directory.
% Compile with: xelatex framework_spec.tex
\usepackage{fontspec}
\defaultfontfeatures{Path = fonts/}
\setmainfont{times.ttf}[
  BoldFont       = timesbd.ttf,
  ItalicFont     = timesi.ttf,
  BoldItalicFont = timesbi.ttf
]

% Math font setup with mathspec.
% IMPORTANT: The sasgts* fonts contain only math symbols (Greek, operators),
% NOT the Latin alphabet. So we map Latin letters in math to the regular Times
% italic font, and use sasgts* only for Greek letters.
\usepackage{mathspec}
\setmathsfont(Latin,Digits)[
  Path           = fonts/,
  UprightFont    = times.ttf,
  ItalicFont     = timesi.ttf,
  BoldFont       = timesbd.ttf,
  BoldItalicFont = timesbi.ttf
]{times.ttf}
\setmathsfont(Greek)[
  Path           = fonts/,
  UprightFont    = sasgtsr.ttf,
  ItalicFont     = sasgtsi.ttf,
  BoldFont       = sasgtsb.ttf,
  BoldItalicFont = sasgtsbi.ttf
]{sasgtsi.ttf}

% FALLBACK if the above math font setup fails:
% Comment out the mathspec block above and uncomment the line below.
% \usepackage{mathptmx}  % Times-compatible math symbols using TeX's built-in fonts

% Load hyperref last for best compatibility.
\usepackage{hyperref}
\hypersetup{
  colorlinks=true,
  linkcolor=blue!50!black,
  urlcolor=blue!50!black,
  citecolor=blue!50!black
}

\title{A Prior-Informed Bootstrap Framework for\\Network Credible Intervals Under Non-Random Missingness}
\author{Jonathan H. Morgan, Ph.D.\\
\small Specification draft for developer review\\
\small Generalizing methods from Daniele Bellutta, \emph{Forensic Analysis of Observed Social Networks}}
\date{\today}

\newcommand{\Gobs}{G_{\text{obs}}}
\newcommand{\Gaug}{G_{\text{aug}}}
\newcommand{\Gb}{G_b}
\newcommand{\pinode}{\pi_{\text{node}}}
\newcommand{\piedge}{\pi_{\text{edge}}}
\newcommand{\pitotal}{\pi_{\text{total}}}
\newcommand{\Nadd}{N_{\text{add}}}
\newcommand{\Wadd}{W_{\text{add}}}
\newcommand{\Waug}{W_{\text{aug}}}
\newcommand{\EI}{\text{E/I}}
\newcommand{\Rmat}{R}
\newcommand{\Pmat}{P}

\begin{document}

\maketitle

\begin{abstract}
This document specifies a bootstrap-based framework for constructing approximate Bayesian credible intervals on network metrics computed from incomplete network observations. The framework extends and addresses limitations of the bootstrap methods in Chapters~3 and~4 of Bellutta's dissertation (\emph{Forensic Analysis of Observed Social Networks}) through four innovations: (1)~explicit modeling of node-level missingness governed by a user-supplied prior on the correlation between centrality and missingness; (2)~structural-role binning via the E/I~index, which separates brokers from clustered hubs of equal degree; (3)~rank-conditional probabilistic imputation of edges involving nominated non-respondents, generalizing the global reciprocity rate of \citet{SmithMorganMoody2022}---the author's prior work---to a rank-rank conditional reciprocity matrix; and (4)~an estimate-based (\emph{deficit}) weight-reconstruction step that inverts the assumed removal model, restoring each edge toward its estimated pre-removal weight rather than redistributing recovered weight in proportion to the degraded observation. Together these extensions let the framework handle unweighted networks (which the prior framework could not), networks with distance-semantic weights (via a transformation step), and networks whose structure departs from the stylized models of earlier evaluations. The specification is given in full as the basis for the Julia implementation and its empirical validation.
\end{abstract}

\section{Introduction}

Network analysts routinely work with incomplete observations of an underlying true network. Whether the data come from surveys with non-response, surveillance with limited resources, social media APIs with rate limits, or informant testimony about covert organizations, the observed network~$\Gobs$ is typically only a partial view of some unobserved true network~$T$. Quantifying uncertainty about network measures computed on~$\Gobs$ has long been a difficult problem, particularly when the missingness mechanism is non-random and varies systematically with node properties such as centrality.

This framework provides a principled approach to constructing credible intervals on network metrics in this setting. The user supplies:
\begin{itemize}[leftmargin=*, itemsep=2pt]
  \item An observed network~$\Gobs$
  \item A prior~$\rho \in (-1, 1)$ on the correlation between node centrality and probability of being missing from observation
  \item Their best estimate of the magnitude of missingness: the node fraction~$\pinode$ and the edge-weight fraction~$\piedge$. Because analysts usually have firmer intuition about missing nodes than missing weight, a convenience interface also accepts a single overall figure~$\pitotal$, treated as~$\pinode$ with~$\piedge$ left to its implied floor.
  \item The network measure(s) of interest
\end{itemize}
The framework then generates a posterior sample of plausible true networks consistent with these inputs, computes the requested measures on each, and returns credible intervals derived from the posterior sample.

The motivating use cases include covert network analysis (where central actors invest heavily in concealment), survey-based network studies (where non-response often correlates with status or visibility), social media analysis (where streaming APIs systematically over-represent active users), and any setting where the analyst has domain knowledge about the observation process but cannot fully characterize it. In each case, the framework's prior~$\rho$ allows the analyst to encode this knowledge directly, and the resulting credible intervals reflect both the magnitude and the structure of the assumed missingness.

\section{Relationship to Bellutta's Prior Work}

The framework builds on the bootstrap methods described in Chapters~3 and~4 of Bellutta's dissertation, \emph{Forensic Analysis of Observed Social Networks}, but extends them in several important ways. Understanding these extensions requires briefly recapping the prior framework's scope and limitations.

\subsection{Chapter 3: The Underlying Noise Model}

Bellutta's Chapter~3 introduced a noise model used to simulate missing data on known true networks for the purpose of evaluating link-inference algorithms. The mechanism removes integer edge-weight units uniformly at random across all weight units in the network~--- equivalent to sampling edges with probability proportional to current weight and decrementing by one, with replacement. The total amount removed is a specified percentage of total edge weight.

This mechanism has several important properties:
\begin{itemize}[leftmargin=*, itemsep=2pt]
  \item It is MCAR at the level of weight units, not at the level of edges
  \item High-weight edges lose more weight in expectation but rarely vanish
  \item No new edges are created during the removal process
  \item The node set is preserved entirely~--- only edges and weights change
  \item It assumes the network is weighted; the mechanism has no analog for unweighted networks
\end{itemize}

\subsection{Chapter 4: The Bootstrap Framework}

Bellutta's Chapter~4 used Chapter~3's noise model as the basis for a bootstrap approach to constructing confidence intervals. Given an observed incomplete network, five methods were evaluated for generating replicate networks: random augmentation (adding integer weight units uniformly at random), the random dot product graph (RDPG) bootstrap, the masked RDPG bootstrap, the structural perturbation method (SPM), and a non-negative matrix factorization perturbation method (NMFP). Confidence intervals were constructed by computing the metric of interest on each replicate and forming a normal-approximation interval centered at the observed metric value.

The Chapter~4 framework has several limitations that motivate this work:

\begin{enumerate}[leftmargin=*, itemsep=2pt]
  \item \textbf{It assumes MCAR.} All five bootstrap methods operate under the assumption that missingness is uniform across edge-weight units. None can encode prior beliefs about which edges or nodes are more likely to be missing.

  \item \textbf{It cannot handle unweighted networks.} Because Chapter~3's noise model removes integer weight units, the framework cannot operate on networks where edges are simply present or absent. An entire class of common networks falls outside its scope.

  \item \textbf{It cannot handle distance-semantic weights.} The noise model assumes weights represent counts where higher equals stronger. Distance networks, where higher weight means weaker tie, would be processed incorrectly.

  \item \textbf{It does not model node-level missingness.} Chapter~3 removes weight from edges but preserves the node set. In practice, much network missingness is at the node level: non-respondents in surveys, unobserved actors in covert networks, accounts excluded from a social media sample. The Chapter~4 framework has no mechanism for this.

  \item \textbf{It was evaluated only on stylized networks.} The framework's empirical validation relied on uniformly random, small-world, and scale-free network models. Real networks often combine structural features in ways these models do not capture, and the framework's behavior on such networks is unclear.

  \item \textbf{Empirical results showed intervals were often miscalibrated.} Chapter~4's coverage results showed that the bootstrap methods produced intervals that were either too narrow (insufficient coverage) or too wide to be useful. Improvements in both the underlying missingness model and the structural assumptions are warranted.
\end{enumerate}

\subsection{What This Framework Adds}

The framework introduced in this document addresses each limitation:

\paragraph{Explicit node-level missingness.} The framework treats missing nodes as a first-class mechanism, separate from edge-weight loss. The parameter $\pinode$ governs the fraction of nodes believed missing, and a synthetic-node augmentation step adds plausible missing nodes to each replicate. Edge-weight missingness ($\piedge$) becomes a secondary, floored dial: node additions and~$\rho$ fix its lower bound, and any weight beyond that floor is allocated by Stage~2 across observed and imputed edges, targeting each edge's estimated pre-removal deficit rather than its current (degraded) weight.

\paragraph{Centrality-correlated missingness via a user-supplied prior.} The parameter~$\rho \in (-1, 1)$ encodes the user's domain-knowledge belief about the association between node centrality and probability of being missing from observation. Positive~$\rho$ corresponds to settings where central actors are more likely to be hidden (covert networks, non-response among prominent people); negative~$\rho$ corresponds to settings where central actors are more visible (social media, surveillance prioritizing known figures); zero~$\rho$ recovers MCAR behavior. The framework converts~$\rho$ into a bin-assignment distribution that biases the centrality of added nodes accordingly.

Internally, $\rho$ is interpreted as Kendall's $\tau_b$, the rank-based correlation between the missingness indicator (binary) and node centrality (continuous). This choice matches the user's rank-order intuition about missingness (``low-ranked nodes are likely to be missing'' is a statement about ranks, not raw centrality values), and it avoids an asymmetric-ceiling problem that arose under Pearson correlation in early framework development. On heavy-tailed networks (e.g., Marvel co-appearance, citation networks), Pearson correlation between missingness and centrality saturates at a much lower magnitude on the negative side than the positive side, because the many low-centrality nodes have near-identical values that fail to discriminate missing from observed. Kendall $\tau_b$ on ranks removes this sign asymmetry, giving comparable ceilings across networks and across the sign of~$\rho$. Two milder limits remain~--- an absolute ceiling that tightens as the missingness rate falls, and slight residual asymmetry at the lowest rates on the most heavy-tailed networks~--- but neither reintroduces the Pearson pathology. The user-facing parameter $\rho \in (-1, 1)$ is unchanged from a standard correlation interface; the metric switch is internal.

\paragraph{Native support for unweighted networks.} Because node-level missingness is independent of edge-weight semantics, the framework operates on unweighted networks by skipping the weight-addition (Stage~2) step entirely. This is a substantive extension: the Chapter~4 framework cannot handle unweighted networks at all, but unweighted networks are central to several lines of prior work on missing-data bias in survey-based networks \citep[e.g.,][of which the present author was a co-author]{SmithMorganMoody2022}, and supporting them is one of the primary motivations for treating node missingness explicitly.

\paragraph{Distance-semantic weight transformation.} The framework provides automatic transformation of distance-weighted networks to a count-like scale before processing, with three transformation methods (scaled reciprocal, max-minus-distance, exponential decay). Credible intervals returned by the framework apply to the transformed network, with explicit caveats about back-transformation.

\paragraph{Reduced reliance on stylized network models.} Rather than fitting a parametric generative model (e.g., RDPG) to the observed network and bootstrapping from it, the framework uses empirical rank-rank connectivity directly. Each replicate's added-node edges are sampled from the observed network's empirical structure, conditional on the centrality and brokerage bins involved. The framework therefore makes no assumption that the true network follows any particular stylized form~--- it only assumes that the rank-rank structure observed in~$\Gobs$ is representative of the true network's rank-rank structure.

\paragraph{Structural-role binning via the E/I index.} The framework uses two binning dimensions by default: degree centrality rank and the E/I~index. The E/I~index, derived from CHAMP community detection on~$\Gobs$ (a resolution sweep over the Leiden algorithm), distinguishes brokers (whose ties cross community boundaries) from clustered hubs (whose ties remain within a single community). This addresses a key weakness of degree-only binning: two nodes with identical degree may play fundamentally different structural roles, and treating them as exchangeable miscalibrates credible intervals for measures sensitive to bridging behavior (e.g., betweenness centrality). The framework falls back to degree-only binning automatically when community structure is not meaningful.

\paragraph{Rank-conditional probabilistic imputation for nominated non-respondents.} In survey, surveillance, and informant-based settings, the analyst frequently knows that certain unobserved nodes exist because they appear as targets of edges from observed (respondent) nodes. These ``nominated non-respondents'' carry partial information~--- their incoming edges are known~--- but their outgoing behavior must be imputed. \citet{SmithMorganMoody2022}, the present author's prior work on missing-data bias in survey-based networks, found that probabilistic imputation, using the global reciprocity rate to impute reverse-direction edges, worked well across many directed-network settings. The framework generalizes this approach by computing a rank-conditional reciprocity matrix~$\Rmat$ where each cell gives the reciprocity rate for forward edges between specific rank-rank cell pairs, rather than a single global rate. This captures structural variation in reciprocity behavior (e.g., peripheral nodes reciprocate at different rates than central hubs, brokers exhibit characteristic asymmetries) while remaining computationally tractable, in contrast to model-based imputation approaches (e.g., ERGMs) that scale poorly. The framework distinguishes nominated non-respondents (handled via~$\Rmat$-based imputation in a new Stage~0.5) from non-nominated non-respondents (handled via the rank-rank matrix~$\Pmat$ in Stage~1).

\paragraph{Bayesian framing.} The framework's output is an approximate posterior sample over plausible true networks, conditional on the observed network and the user's prior. Credible intervals are derived directly from quantiles of the posterior sample, without the normal-approximation assumption of Chapter~4. The framework is not derived from an explicit likelihood and prior density formulation, so its output is "approximately" Bayesian in spirit rather than strictly so, but the conceptual framing is fundamentally different from Chapter~4's frequentist confidence intervals.

\section{Algorithm Specification}

This section provides the complete algorithmic specification of the framework. The algorithm proceeds in three phases: a setup phase executed once on the observed network, a replicate generation phase executed~$B$ times to produce the posterior sample, and an aggregation phase that summarizes the posterior.

\subsection{Inputs}

\paragraph{Required inputs.}
\begin{itemize}[leftmargin=*, itemsep=2pt]
  \item $\Gobs$: an observed network with $N$ nodes. The network may be directed or undirected, weighted (with non-negative integer edge weights representing counts) or unweighted. Distance-semantic networks must be transformed (see Section~\ref{sec:weight-semantics}).
  \item $\rho \in (-1, 1)$: the prior correlation between node centrality and probability of being missing.
  \item One or more network measures $m$ for which credible intervals are desired.
\end{itemize}

\paragraph{Missingness specification.} The user provides either the simple or advanced interface:
\begin{itemize}[leftmargin=*, itemsep=2pt]
  \item Simple: $\pitotal \in [0, 1)$, the overall fraction of the true network believed unobserved. The framework defaults this to $\pinode = \pitotal$, $\piedge = 0$.
  \item Advanced: $\pinode \in [0, 1)$ and/or $\piedge \in [0, 1)$, the node-level and edge-level missingness fractions respectively. If only one is supplied, the other defaults to~0.
\end{itemize}
At least one of $\pinode$ or $\piedge$ must be positive. For unweighted networks, $\piedge$ is structurally inapplicable and must be~0; $\pinode > 0$ is required.

\paragraph{Sign convention for $\rho$.}
\begin{itemize}[leftmargin=*, itemsep=2pt]
  \item $\rho > 0$: central nodes are \emph{more} likely to be missing (e.g., they hide successfully, refuse to respond, are hard to reach)
  \item $\rho < 0$: central nodes are \emph{less} likely to be missing (e.g., they are prioritized for surveillance, naturally over-captured by streaming APIs)
  \item $\rho = 0$: no correlation between centrality and missingness (MCAR)
\end{itemize}
A negative~$\rho$ does not denote a negative probability; it denotes a negative correlation between centrality and the indicator of being missing.

The framework interprets $\rho$ internally as Kendall's $\tau_b$, the rank-based correlation between centrality and the missingness indicator. The user-facing parameter is a standard correlation in $(-1, 1)$; the rank-based interpretation is internal and not exposed to the user except in diagnostic output. See Section~\ref{sec:feasibility} for the feasibility constraints that determine whether a given~$\rho$ is structurally achievable on a given network at a given missingness rate.

\paragraph{Guidance table for $\rho$.} Table~\ref{tab:rho-guidance} provides illustrative values of~$\rho$ for common observation processes. Users should set~$\rho$ based on domain knowledge, not by estimating it from~$\Gobs$.

\begin{table}[h]
\centering
\small
\begin{tabular}{p{5.5cm} p{3.5cm} p{1.8cm} p{1.8cm}}
\toprule
\textbf{Observation process} & \textbf{Direction} & \textbf{Typical $|\rho|$} & \textbf{Signed $\rho$} \\
\midrule
Complete-roster survey, high response rate & No bias & 0.0 & 0.0 \\
Social media via streaming/firehose API & Central more visible & 0.3 -- 0.5 & $-0.3$ to $-0.5$ \\
Snowball sample from initial respondents & Central more visible & 0.4 -- 0.6 & $-0.4$ to $-0.6$ \\
Survey with refusals among prominent people & Central more hidden & 0.3 -- 0.5 & $+0.3$ to $+0.5$ \\
Covert network from informant testimony & Central more hidden & 0.4 -- 0.7 & $+0.4$ to $+0.7$ \\
Covert network from police surveillance & Central more visible & 0.5 -- 0.7 & $-0.5$ to $-0.7$ \\
Mixed police data & Variable & Variable & User judgment \\
\bottomrule
\end{tabular}
\caption{Guidance for setting~$\rho$ based on observation process. Values are illustrative starting points, not prescriptions. Whether a chosen $\rho$ is structurally achievable on a given network at a given missingness rate is a separate question; the framework surfaces this via \texttt{feasible\_rho\_range} (Section~\ref{sec:feasibility}). On heavy-tailed networks at low missingness rates, the rate-bounded ceiling $\sqrt{2 p (1-p)}$ may limit the achievable $|\rho|$ below the values in this table; the framework will then operate at the practical ceiling and report both nominal and realized values.}
\label{tab:rho-guidance}
\end{table}

\paragraph{Optional inputs.}
\begin{itemize}[leftmargin=*, itemsep=2pt]
  \item \texttt{community\_labels}: user-supplied community assignments, overriding CHAMP detection
  \item \texttt{weight\_semantics}: \texttt{"count"} (default) or \texttt{"distance"}; selects transformation behavior
  \item \texttt{partially\_observed\_nodes}: a list of node identifiers for nodes that appear in~$\Gobs$ as edge targets (i.e., they have incoming edges from respondents) but did not respond themselves. These are the ``nominated non-respondents'' handled by Stage~0.5. If not supplied, all node-level missingness is assumed to involve non-nominated nodes (handled by Stage~1).
  \item \texttt{allocation}: the Stage~2 weight-allocation rule. \texttt{:deficit} (default) inverts the removal model, restoring each edge toward its estimated pre-removal weight; \texttt{:observed} reproduces the earlier current-weight-proportional rule (the Bellutta MCAR corner). The two coincide at $\rho = 0$, and the choice has no effect on unweighted networks, where Stage~2 is skipped.
\end{itemize}

\paragraph{Algorithm hyperparameters with defaults.}
\begin{itemize}[leftmargin=*, itemsep=2pt]
  \item $B$: number of posterior samples (default 1000)
  \item $K$: number of degree-rank bins. Default is automatic selection (\texttt{K = :auto}) via the function \texttt{find\_optimal\_K}, which iterates $K$ from a maximum tractable value down to a floor and returns the largest $K$ that satisfies the framework's algorithmic requirements (see Section~\ref{sec:three-priors}). Larger $K$ produces a Kendall $\tau_b$ ceiling closer to the rate-bounded absolute (fewer bin-level ties on the centrality side), at the cost of fewer observed nodes per bin. The default $K_{\max}$ is $\min(20, \lfloor N_{\text{obs}} / 5 \rfloor)$, the default $K_{\min}$ is~4. Users may override with an explicit integer when manual control is desired.
  \item $J$: number of E/I bins (default 3, with semantic categories internal-leaning [$\EI \leq -0.33$], balanced [$-0.33 < \EI < 0.33$], broker-leaning [$\EI \geq 0.33$])
\end{itemize}

\subsection{Feasibility of $\rho$}
\label{sec:feasibility}

Not every $\rho \in (-1, 1)$ is structurally achievable on a given network at a given missingness rate. Two distinct ceilings constrain the realized correlation:

\paragraph{Rate-bounded absolute ceiling.} Under Kendall's $\tau_b$ between a binary indicator (missingness, with proportion $p$) and a continuous variable (centrality), the realized correlation is structurally bounded by
\[
|\tau_b|_{\max} = \sqrt{\frac{2 \cdot k \cdot (n - k)}{n \cdot (n - 1)}} \approx \sqrt{2 \cdot p \cdot (1 - p)}
\]
where $k$ is the number of missing nodes and $p = k/n$. This bound depends only on the missingness rate, not on the network structure. At $p = 0.10$, $|\tau_b|_{\max} \approx 0.42$; at $p = 0.25$, $|\tau_b|_{\max} \approx 0.61$; at $p = 0.50$, $|\tau_b|_{\max} \approx 0.71$. The framework cannot produce realized $\rho$ exceeding this absolute ceiling no matter how the sampling is configured.

\paragraph{Practical tie-bounded ceiling.} On networks where centrality has many ties (most unweighted networks, where multiple nodes share the same degree), the achievable $|\tau_b|$ is further reduced below the absolute ceiling. Centrality ties contribute to the denominator of Kendall $\tau_b$, lowering the maximum that any sampling configuration can reach. The practical ceiling is network-specific and depends on the centrality distribution's tie structure.

\paragraph{Feasibility check.} The framework provides a public function \texttt{feasible\_rho\_range} that probes the sampling mechanism at extreme $\beta$ configurations to estimate the practical $[\rho_{\min}, \rho_{\max}]$ achievable on a given network at a given rate. \texttt{build\_reconstruction\_corpus} calls this function at startup and clamps the user's $\rho$ into the practical range, conditioning the entire sample on the clamped value and recording whether an adjustment occurred.

\paragraph{What happens when $\rho$ exceeds the practical ceiling.} The framework clamps the requested $\rho$ to the feasible envelope and conditions the sample on the clamped value. Because a credible interval is meaningful only relative to the prior actually enforced, the clamped (conditioned) $\rho$ becomes the operative prior, and both the requested and conditioned values are recorded. Within the conditioned target, the Step~5 $\beta$-solver may still return \texttt{:ceiling\_hit} if that value exceeds the $K$-dependent bin ceiling, in which case $\beta$ saturates at the bracket endpoint and the realized $\rho$ is the maximum the binning can produce. Diagnostic output reports nominal (requested), conditioned (clamped), and realized $\rho$, so the user can interpret the credible intervals against the prior that was actually applied.

\paragraph{Asymmetry between $+\rho$ and $-\rho$.} Under Kendall $\tau_b$ with rank-based sampling weights (see Step~5), feasibility ceilings are typically symmetric: $|\rho_{\min}| \approx |\rho_{\max}|$. This is a substantive improvement over Pearson correlation, where heavy-tailed networks produced dramatically asymmetric ceilings (e.g., $\rho_{\max} \approx +0.4$ but $\rho_{\min} \approx -0.03$ on Marvel under Pearson). The \texttt{is\_asymmetric} flag in \texttt{feasible\_rho\_range} output retains a diagnostic for unusual cases but rarely fires under Kendall.

\subsection{Weight Semantics and Transformation}
\label{sec:weight-semantics}

The framework assumes edge weights are non-negative integer counts representing interactions or analogous additive quantities where higher weight indicates a stronger tie. Networks where weights represent distances, dissimilarities, or other "lower = stronger" quantities must be transformed to a count-like scale before processing.

When \texttt{weight\_semantics = "distance"} is set, the framework applies one of the following transformations in a preprocessing step:
\begin{itemize}[leftmargin=*, itemsep=2pt]
  \item \textbf{Scaled reciprocal (default):} $w' = \mathrm{round}(c / w)$ with $c = \max(w)$. Requires distances to be strictly positive; zero-distance edges, if any, take the maximum observed transformed weight.
  \item \textbf{Max-minus-distance:} $w' = \mathrm{round}(\max(w) - w)$. Linear inversion; suitable when distances have a meaningful upper bound. Distance-zero maps to maximum count; maximum distance maps to zero (treated as no edge).
  \item \textbf{Exponential decay:} $w' = \mathrm{round}(c \cdot \exp(-w / \tau))$ for user-supplied decay parameter~$\tau$; $c$ defaults to a value chosen so that the median transformed weight is approximately~10.
\end{itemize}
After transformation, all subsequent steps operate on the transformed network. Credible intervals returned by the framework apply to the transformed network, not the original distances. Users wishing to make inferences about original-scale metrics must back-transform results carefully, recognizing that nonlinear transformations may distort the interval structure. For rank-based measures (e.g., who has highest degree centrality), rank ordering is often preserved across reasonable transformations; for distance-dependent measures (e.g., closeness, betweenness on weighted shortest paths), the transformation fundamentally changes what the measure represents.

\subsection{The Bin-Exchangeability Assumption}

The framework treats all observed nodes within a bin as exchangeable. When sampling edges for added nodes, bin-level connection probabilities and weights are used rather than node-level. This is the framework's central simplifying assumption.

By default, binning is two-dimensional: each node is characterized by both its degree rank and its E/I~index. This recognizes that two nodes with the same degree may play fundamentally different structural roles. Within each (degree-bin~$\times$~E/I-bin) cell, individual differences in connection patterns are still averaged out; to the extent that within-cell structural variation matters for the user's measure of interest, credible intervals may be miscalibrated. The framework is best suited to networks where degree and~$\EI$ together capture the dominant structural patterns; networks dominated by attribute homophily, geographic structure, or community structure orthogonal to degree-and-brokerage remain difficult cases.

\subsection{Structural-Role Binning via the E/I Index}

The framework always uses the E/I~index as a second binning dimension alongside degree. The E/I~index measures how much a node's ties cross community boundaries, distinguishing brokers (high $\EI$) from clustered hubs (low $\EI$) even when they have the same degree.

\paragraph{Definition.} For each node~$i$,
\[
\EI_i = \frac{E_i - I_i}{E_i + I_i},
\]
where $E_i$ is the count (or weight, for weighted networks) of node~$i$'s external ties (crossing community boundaries) and $I_i$ is the count (or weight) of its internal ties. The index ranges from~$-1$ (purely internal) to~$+1$ (purely external).

\paragraph{Community detection.} If user-supplied community labels are not provided, the framework runs CHAMP community detection (a resolution sweep over the Leiden algorithm) on~$\Gobs$. For directed networks, CHAMP uses its directed modularity; for weighted networks, it uses edge weights directly. Community detection runs once at setup and is held fixed across all replicates.

\paragraph{Fallback behavior.} If community detection returns a single community, or if E/I has insufficient variance to support~$J$ bins with at least three nodes per bin, the framework issues a diagnostic message and falls back to degree-only binning for the remainder of the run. This handles edge cases gracefully: networks with no meaningful community structure, very small networks where community detection is unreliable, and networks where all nodes have similar bridging behavior. The fallback is automatic; no user intervention is required. The binning mode used (2D or fallback to 1D) is reported in diagnostics.

\paragraph{Two-dimensional binning.} When E/I binning is active, each observed node is assigned both a degree bin (1~to~$K$) and an E/I bin (1~to~$J$). The rank-rank connectivity matrix is indexed by (degree-bin,~EI-bin) cell pairs. Because per-cell estimates can be noisy when cells contain few nodes, the framework applies Laplace smoothing to the connection probabilities~$P$ in all cases.

\paragraph{Bin assignment for added nodes.} Added nodes are assigned degree-bins via the $\rho$-weighted distribution. Their E/I bins are drawn from the empirical conditional distribution of E/I bin given the assigned degree bin in~$\Gobs$. Central brokers are therefore sampled when central added nodes are drawn, in proportion to how common brokers are at high centrality in the observed network. Users with specific prior beliefs about brokerage-missingness correlation (distinct from centrality-missingness) may extend the framework with a separate $\rho_{\EI}$ parameter; this is left as an advanced extension.

\subsection{Nominated Non-Respondents and the Reciprocity Matrix}
\label{sec:nominated}

In many observation processes, the analyst has explicit knowledge of some unobserved nodes because they appear as targets of edges from observed nodes. In survey settings, these are non-respondents who were nominated by respondents; in surveillance settings, they are actors named in observed activity reports; in social media samples, they are accounts mentioned by sampled users. The framework refers to these as \emph{nominated non-respondents} (or, more generally, partially observed nodes).

The user identifies these nodes through the optional input \texttt{partially\_observed\_nodes}. When supplied, the framework treats them differently from the synthetic ``added'' nodes generated by Stage~1: they are placed in the network using their observed incoming edges, and their outgoing edges are imputed using a rank-conditional reciprocity model.

\paragraph{The reciprocity matrix $\Rmat$ (directed networks only).} For each ordered pair of cells $(c, c')$, define
\[
\Rmat_{c, c'} = \Pr\left( j \to i \text{ exists} \,\middle|\, i \to j \text{ exists}, \; i \in c, \; j \in c' \right).
\]
In words: $\Rmat_{c, c'}$ is the probability that a forward edge from a node in cell~$c$ to a node in cell~$c'$ is reciprocated by a reverse edge from~$c'$ back to~$c$. The matrix is estimated empirically from respondent-respondent dyads in~$\Gobs$ with Laplace smoothing:
\[
\Rmat_{c, c'} = \frac{|\{(i, j) : i \in c, \, j \in c', \, w(i, j) > 0, \, w(j, i) > 0\}| + 1}{|\{(i, j) : i \in c, \, j \in c', \, w(i, j) > 0\}| + 2}.
\]
Unlike the unconditional matrix~$\Pmat$, the reciprocity matrix~$\Rmat$ is fundamentally asymmetric: $\Rmat_{c, c'}$ and $\Rmat_{c', c}$ can be very different in networks with structurally asymmetric tie patterns (e.g., hub-and-spoke or core-periphery networks).

\paragraph{Generalization of probabilistic imputation.} The matrix~$\Rmat$ generalizes the single global reciprocity rate used in the probabilistic imputation of \citet{SmithMorganMoody2022}. Where the prior approach used one value for all dyads (e.g., 0.25 in their illustrative example), the framework estimates a separate value for each rank-rank cell pair. With broker detection enabled, $\Rmat$ also captures characteristic asymmetries between brokers and clustered hubs that a single rate would miss.

\paragraph{Bin assignment for nominated non-respondents.} Each nominated non-respondent is placed in degree and E/I bins based on its \emph{observed} incoming edges. This means the bin assignment may under-estimate the node's true centrality (since outgoing edges are missing), but it uses only information actually available in~$\Gobs$. Adjusting for under-observation is left as an advanced extension.

\paragraph{Undirected networks.} Nominated non-respondents are a directed notion: their incoming ties are observed while their outgoing ties are missing. Undirected networks have no in/out asymmetry, so a missing undirected node is modeled as a full removal (handled by Stage~1), not as a nomination. Accordingly, the framework rejects \texttt{partially\_observed\_nodes} on undirected~$\Gobs$ rather than imputing them. This is a known limitation: undirected networks dominated by edges that respondents failed to report (measurement error) are less well-handled by the framework.

\subsection{Edge Representation Convention}

The framework operates on networks that may be directed or undirected, weighted or unweighted:
\begin{itemize}[leftmargin=*, itemsep=2pt]
  \item Directed networks: centrality is in-degree; rank-rank matrices index ordered pairs; added-node edges are sampled separately for outgoing and incoming directions.
  \item Undirected networks: centrality is total degree; rank-rank matrices are symmetric and index unordered pairs; added-node edges are sampled once per pair.
  \item Weighted networks: each imputed tie enters at weight~1 (the floor), with diffuse additional weight redistributed in Stage~2; added-node tie counts are drawn from Poisson tendencies.
  \item Unweighted networks: Stage~2 is skipped; added-node edges have weight~1 by construction; the rank-rank weight matrix is unnecessary.
\end{itemize}
Self-loops are ignored throughout. Multigraphs must be collapsed to weighted single-edge form before input.

\subsection{Step 0 --- Preprocessing}

\paragraph{Step 0a: Weight transformation (if applicable).}
If \texttt{weight\_semantics = "distance"}, apply the selected transformation to~$\Gobs$. All subsequent steps operate on the transformed network. Transformation parameters are recorded for output metadata.

\paragraph{Step 0b: Resolve missingness parameters.}
If~$\Gobs$ (post-transformation, if applicable) is unweighted, set $\piedge = 0$ regardless of user input, issue a warning if the user supplied a non-zero $\piedge$, and require $\pinode > 0$. Otherwise: if $\pinode$ and/or $\piedge$ are explicitly provided, use them as given (defaulting the unspecified one to~0); else if $\pitotal$ is provided, set $\pinode = \pitotal$ and $\piedge = 0$. Validate that $\pinode, \piedge \in [0, 1)$ and at least one is positive.

\subsection{Setup Phase}

\paragraph{Step 1: Compute observed centrality.}
For directed~$\Gobs$, set centrality as in-degree per node; this is the rho basis against which missingness is correlated, and it is deliberately the destination endpoint only. For undirected~$\Gobs$, set centrality as degree. Rank nodes by centrality; ties broken consistently.

\paragraph{Step 1.5: Detect community structure and compute E/I.}
If user-supplied community labels are provided, use them. Otherwise, apply CHAMP to~$\Gobs$ (directed modularity if directed; using weights if weighted). Check whether the community structure is meaningful: more than one community detected, and E/I values supporting~$J$ bins with at least three nodes each. If meaningful, compute~$\EI_i$ for each observed node and rank by~$\EI$. If not, fall back to degree-only binning for the remainder of the run and log a diagnostic.

\paragraph{Step 2: Bin observed nodes.}
Partition observed nodes into~$K$ equally-sized degree bins, with bin~1 the most peripheral and bin~$K$ the most central. If 2D binning is active, additionally partition into~$J$ E/I bins using the semantic thresholds ($-0.33$, $+0.33$) for $J = 3$, or equally-sized rank bins for other~$J$.

\paragraph{Step 3: Compute the rank-rank connectivity matrix $\Pmat$.}
For each cell pair $(c, c')$ (ordered for directed networks, unordered for undirected), compute the Laplace-smoothed connection probability
\[
p_{c, c'} = \frac{e_{c, c'} + 1}{n_{c, c'} + 2},
\]
where $e_{c, c'}$ is the number of observed edges from cell~$c$ to cell~$c'$ and $n_{c, c'}$ is the number of ordered node pairs between the two cells. For weighted networks, also compute the conditional mean weight $w_{c, c'}$: the mean weight of edges from cell~$c$ to cell~$c'$ given that an edge is present. For unweighted networks, only $\Pmat$ is computed. The matrices are computed over all respondent-respondent dyads in~$\Gobs$ (excluding any partially observed nodes, which are placed using their own observed edges separately).

\subsection{Stage 2: Weight Allocation}

Stage~2 returns the diffuse weight budget~$\Wadd$ to the network. The guiding idea is to undo the Chapter~3 removal rather than to pile weight onto whatever survived it. The removal stripped weight from each edge in proportion to a per-edge hazard set by the $\rho$-field; Stage~2 reads that same field, estimates how much each edge is likely to have lost, and returns weight to the edges with the largest estimated shortfall.

Let $w_e$ be the observed (post-removal) weight of edge~$e$, and let $g_e \in [0, 1]$ be its endpoint-survival tilt~--- the edge's exposure to removal, equal to one minus the probability that both of its endpoints escaped removal, with each endpoint's survival probability set by the $\rho$-field. The framework models the surviving fraction of an edge as $\exp(-\lambda g_e)$ for a single rate~$\lambda$ shared across all edges, so the estimated pre-removal weight of~$e$ is $w_e \exp(\lambda g_e)$ and its estimated deficit is
\[
  d_e = w_e \left( \exp(\lambda g_e) - 1 \right).
\]
The rate~$\lambda$ is fixed by the budget: it is the unique value for which the deficits sum to the weight to be added,
\[
  \sum_e w_e \left( \exp(\lambda g_e) - 1 \right) = \Wadd .
\]
The left-hand side rises monotonically in~$\lambda$ from zero, so a bisection finds the root in a few iterations ($\lambda$ is capped at a large finite value to guard against degenerate inputs). Stage~2 then draws the~$\Wadd$ units multinomially with probabilities proportional to~$d_e$, restoring each edge to its estimated true weight~$w_e \exp(\lambda g_e)$ in expectation.

Two properties follow directly. When $\rho = 0$ the tilt is flat~--- $g_e$ is the same for every edge~--- so the common factor~$\exp(\lambda g_e)$ cancels and the rule reduces exactly to current-weight-proportional allocation, the Bellutta MCAR corner recovered by \texttt{allocation = :observed}. And an edge driven to zero ($w_e = 0$) has zero deficit and receives nothing: Stage~2 redistributes weight among edges already present in~$\Gobs$ but never resurrects a vanished one, in keeping with the framework's assumption that edge presence among respondents is observed.

\section{Key Assumptions}

The framework makes the following assumptions, which users should consider when interpreting results:

\begin{enumerate}[leftmargin=*, itemsep=4pt]
\item \textbf{Edge presence among respondent nodes is fully observed.} Respondents report their own ties accurately: for any two respondent nodes (observed nodes that are not partially observed nominated non-respondents), the recorded presence or absence of an edge is taken as correct, and the framework never creates or deletes an edge between two respondents. For weighted networks, Stage~2 may redistribute weight onto existing edges, but it adds no respondent-respondent edge; for unweighted networks, no respondent-respondent edge is modified at all. The assumption covers respondent-respondent dyads only. A nominated non-respondent is observed solely through the ties respondents reported \emph{to} it~--- its own outgoing ties are unobserved and are imputed in Stage~0.5~--- and non-nominated added nodes are absent entirely and are introduced in Stage~1; neither is assumed fully observed.

  \item \textbf{Rank-rank connection structure in~$\Gobs$ is representative of the true network.} The probability that a node in cell~$c$ connects to a node in cell~$c'$ in the truth equals the empirical $p_{c,c'}$ from~$\Gobs$ in expectation.

  \item \textbf{Rank-conditional reciprocity behavior is representative.} For directed networks, the empirical reciprocity rate $\Rmat_{c, c'}$ observed among respondent-respondent dyads is representative of reciprocity behavior involving nominated non-respondents at the same rank-rank cell pair. This generalizes the assumption in \citet{SmithMorganMoody2022} that a single global reciprocity rate could be used for imputation; the framework's assumption is that reciprocity varies by rank-rank cell but is otherwise constant.

  \item \textbf{$\rho$ is a genuine prior, not estimated from~$\Gobs$.} Users supply~$\rho$ based on domain knowledge about the observation process, not by looking at the data. The framework interprets~$\rho$ internally as Kendall's $\tau_b$ between the missingness indicator and node centrality. When the user's $\rho$ exceeds the practical feasibility ceiling (Section~\ref{sec:feasibility}), the framework clamps it to the feasible envelope, conditions the sample on the clamped value, and reports nominal, conditioned, and realized $\rho$ via diagnostics; the contract is algorithmic compliance with what is structurally achievable, not silent overrides of the user's request.

  \item \textbf{Bin membership, not node identity, governs sampling.} All observed nodes in the same cell are treated as exchangeable. This is the framework's central simplifying assumption.

  \item \textbf{Observed-node properties are fixed across replicates.} Degree ranks, E/I ranks, and community labels are computed once on~$\Gobs$ and held constant. Added nodes are assigned bins but do not enter the rankings that define bins. Nominated non-respondents are binned based on their observed incoming edges only.

  \item \textbf{Edge weights (when present) are integer-valued counts of interactions or analogous higher-is-stronger additive quantities.} Distance, dissimilarity, or lower-is-stronger semantics are not supported directly; users must specify a transformation method, after which credible intervals apply to the transformed scale.

  \item \textbf{Default behavior treats missingness as node-level.} Edge-level error on observed nodes ($\piedge$) is opt-in via the advanced interface or non-zero explicit specification, and is unavailable for unweighted networks.

  \item \textbf{Community structure derived from~$\Gobs$ reflects the true community structure.} The detected community structure---CHAMP output, or user-supplied labels---is treated as fixed and correct. To the extent the true community structure differs from what is detected in~$\Gobs$~--- particularly because missingness may obscure community boundaries~--- the E/I~index becomes a noisier signal of bridging behavior.

  \item \textbf{Higher-order structural patterns beyond degree and E/I are not directly modeled.} Triadic closure, attribute homophily, and other patterns appear in replicate networks only to the extent that they correlate with the binning variables. Networks dominated by patterns orthogonal to the binning are less faithfully reproduced.

  \item \textbf{The output is an approximate Bayesian posterior.} The framework produces samples from a distribution consistent with the user's prior, but does not derive from an explicit likelihood and prior density formulation.
\end{enumerate}

\section{What the Framework Does Not Model}

The following are explicitly outside the scope of this framework:

\begin{itemize}[leftmargin=*, itemsep=2pt]
\item Creation of edges between two respondent nodes (the framework assumes edge \emph{presence} among respondents is correctly observed; only \emph{weights} are uncertain, and only for weighted networks). Edges involving nominated non-respondents \emph{are} created/imputed by Stage~0.5.
\item Respondent measurement error: cases where a respondent failed to report an edge to another respondent that actually exists. The framework treats all respondents as accurate reporters of their own edges.
\item A user-specified non-uniform edge-weight missingness mechanism. The structure of edge-weight loss is governed by the $\rho$-field, not by a separate dial. Stage~2 inverts the Chapter~3 removal model: from the endpoint-survival tilt the $\rho$-field assigns to each edge, it estimates the edge's pre-removal weight and allocates the budget~$\Wadd$ toward the resulting per-edge deficits, so each edge is restored to its estimated true weight in expectation (the \emph{deficit} allocation, used by default). The earlier rule, which distributes~$\Wadd$ in proportion to current weight times that tilt, is retained as an option (\texttt{allocation = :observed}); it reproduces the current-weight-proportional behavior and coincides with the deficit allocation at $\rho = 0$ (the Bellutta MCAR corner), where the tilt is flat. Neither rule exposes a separate dial for edge-weight non-uniformity.
\item Spurious edges in the observed network (the framework assumes observed structure is correct, not over-reported)
\item Higher-order structural patterns beyond degree and~$\EI$ (triadic closure, attribute homophily, community structure orthogonal to the bins are captured only incidentally)
\item Changes in node ranks or community memberships due to the augmentation process (all are fixed at~$\Gobs$)
\item Self-loops in added-node sampling (excluded by convention)
\item Multigraphs in input (must be collapsed to weighted single-edge form before running)
\item Edge-weight missingness for unweighted networks (structurally impossible)
\item Native handling of distance, dissimilarity, or lower-is-stronger weight semantics (must be transformed to count-like scale via the framework's preprocessing step)
\item Separate prior on brokerage-missingness correlation ($\rho$ governs degree-missingness; E/I assignment is conditional on degree)
\item Uncertainty in the detected community structure (CHAMP output, or user-supplied labels, is treated as fixed)
\item Reciprocity-based imputation for undirected networks. The reciprocity matrix~$\Rmat$ is undefined for undirected input~--- every undirected edge is mutual by construction~--- and Stage~0.5 imputes a non-respondent's missing \emph{outgoing} ties, which has no undirected analogue. The framework therefore rejects undirected input carrying nominated non-respondents rather than imputing them.
\item Silent enforcement of structurally infeasible $\rho$ priors. When the user's $\rho$ exceeds the rate-bounded or tie-bounded ceiling, the framework clamps it to the feasible envelope, conditions the sample on the clamped value, and emits a diagnostic reporting nominal, conditioned, and realized $\rho$. The clamping is explicit and recorded; the framework does not silently substitute a feasible value for the request or proceed as though the request had been honored.
\end{itemize}

\section{Next Steps}

This document specifies the algorithm as the basis for its Julia implementation and empirical validation. The work ahead spans completing the user-facing layer, building an evaluation harness for the known-ground-truth case, and carrying out the validation program:
\begin{enumerate}[leftmargin=*, itemsep=2pt]
  \item Completing the orchestration layer that draws the~$B$ replicates from the reconstruction procedure and assembles the requested measures into a credible interval~--- the single entry point that turns an observed network and a prior into an interval.
  \item Implementing the calibration diagnostics for the known-ground-truth setting: evaluating each measure at the true~$\rho$ as a best-possible-world reference and at bracketing assumed-$\rho$ values, then scoring how far the off-truth intervals depart from the truth-conditioned one.
  \item Building the distance-semantic weight-transformation module (Section~\ref{sec:weight-semantics}), which maps lower-is-stronger weights onto the count-like scale the framework requires.
  \item Empirical validation on a designed corpus of networks with known ground truth, spanning weight and degree heterogeneity under both edge- and node-level degradation across a range of~$\rho$.
  \item Validation on real network data, including the covert and survey-based networks used in prior work.
  \item Comparison against the Chapter~4 bootstrap methods, to assess whether the framework's intervals are better calibrated under non-random missingness.
  \item Extension to multidimensional (multiplex and/or dynamic) networks, building on the tensor representations developed in Chapter~3.
\end{enumerate}

\bibliographystyle{plainnat}
\bibliography{references}

\end{document}